\section{AIME I 1998}
        \begin{enumerate}
        	\item (\textit{AIME I 1998}) For how many values of $k$ is $12^{12}$ the \href{https://artofproblemsolving.com/wiki/index.php?title=Least\_common\_multiple}{least common multiple} of the positive integers $6^6$, $8^8$, and $k$?


 
 \textbf{Answer: }025



\textbf{Solution 1}

It is evident that $k$ has only 2s and 3s in its prime factorization, or $k = 2^a3^b$. 

 
                \begin{itemize}
                    \item $6^6 = 2^6\cdot3^6$\item $8^8 = 2^{24}$\item $12^{12} = 2^{24}\cdot3^{12}$
                \end{itemize}


                The \href{https://artofproblemsolving.com/wiki/index.php?title=LCM}{LCM} of any numbers an be found by writing out their factorizations and taking the greatest power for each factor. $[6^6,8^8] = 2^{24}3^6$. Therefore $12^{12} = 2^{24}\cdot3^{12} = [2^{24}3^6,2^a3^b] = 2^{\max(24,a)}3^{\max(6,b)}$, and $b = 12$.  Since $0 \le a \le 24$, there are $\boxed{25}$ values of $k$.

 





\textbf{Solution 2}

We want the number of $k$ such that $\operatorname{lcm}(6^6,8^8,k)=12^{12}$.

 Using $\operatorname{lcm}$ properties, this is $\operatorname{lcm}(\operatorname{lcm}(2^6\cdot 3^6,2^{24}),k)=2^{24} \cdot 3^{12}$, or $\operatorname{lcm}(2^{24}\cdot 3^6,k)=2^{24}\cdot 3^{12}$.

 At this point, we realize that $k=2^a\cdot3^b$, as any other prime factors would be included in the $\operatorname{lcm}$. 

 Also, $b=12$ (or the power of $3$ in the $\operatorname{lcm}$ wouldn't be $12$) and $0\le a\le 24$ (or the power of $2$ in the $\operatorname{lcm}$ would be $a$ and not $24$).

 Therefore, $a$ can be any integer from $0$ to $24$, for a total of $25$ values of $a$ and $\boxed{025}$ values of $k$. 

 




	\item (\textit{AIME I 1998}) Find the number of \href{https://artofproblemsolving.com/wiki/index.php?title=Ordered\_pair}{ordered pairs} $(x,y)$ of positive integers that satisfy $x \le 2y \le 60$ and $y \le 2x \le 60$.


 
 \textbf{Answer: }480



\textbf{Solution 1}


                    \begin{figure}[H]
                        \centering
                        \includegraphics[width=0.5\textwidth]{Solution/AIME_I_1998_2_Solution_Figure_1.png}
                    \end{figure}
                    \href{https://artofproblemsolving.com/wiki/index.php?title=Pick\%27s\_theorem}{Pick's theorem} states that:

 $A = I + \frac B2 - 1$The conditions give us four \href{https://artofproblemsolving.com/wiki/index.php?title=Inequality}{inequalities}: $x \le 30$, $y\le 30$, $x \le 2y$, $y \le 2x$. These create a \href{https://artofproblemsolving.com/wiki/index.php?title=Quadrilateral}{quadrilateral}, whose area is $\frac 12$ of the 30 by 30 \href{https://artofproblemsolving.com/wiki/index.php?title=Square}{square} it is in. A simple way to see this is to note that the two triangles outside of the quadrilateral form half of the area of the 30 by 30 square.

 So $A = \frac 12 \cdot 30^2 = 450$. $B$ we can calculate by just counting. Ignoring the vertices, the top and right sides have 14 \href{https://artofproblemsolving.com/wiki/index.php?title=Lattice\_point}{lattice points}, and the two diagonals each have 14 lattice points (for the top diagonal, every value of $x$ corresponds with an integer value of $y$ as $y = 2x$ and vice versa for the bottom, so and there are 14 values for x not counting vertices). Adding the four vertices, there are 60 points on the borders.

 $450 = I + \frac {60}2 - 1$
$I = 421$Since the inequalities also include the equals case, we include the boundaries, which gives us $421 + 60 = 481$ ordered pairs. However, the question asks us for positive integers, so $(0,0)$ doesn't count; hence, the answer is $480$.

 





\textbf{Solution 2}

First, note that all pairs of the form $(a,a)$, $1\le a\le30$ work.

 Now, considered the ordered pairs with $x < y$, so that $x < 2y$ is automatically satisfied. Since $x < y\le 2x$, there are $2x - x = x$ possible values of $y$. Hence, given $x$, there are $x$ values of possible $y$ for which $x < y$ and the above conditions are satisfied. But $2y\le60$, so this only works for $x\le15$. Thus, there are

 $\sum_{i=1}^{15} i=\frac{(15)(16)}{2}$

 ordered pairs. For $x > 15$, $y$ must follow $x < y\le 30$. Hence, there are $30 - x$ possibilities for $y$, and there are

 $\sum_{i=16}^{30}(30-i)=\sum_{i=0}^{14}i=\frac{(14)(15)}{2}$

 ordered pairs.

 By symmetry, there are also $\frac {(15)(16)}{2} + \frac {(14)(15)}{2}$ ordered pairs with $x > y$ and the above criteria satisfied.

 Hence, the total is

 $\frac{(15)(16)}{2}+\frac{(14)(15)}{2}+\frac{(15)(16)}{2}+\frac{(14)(15)}{2}+30=480.$

 





\textbf{Solution 3}

$y\le2x\le60$

 Multiplying both sides by 2 yields:

 $2y\le4x\le120$

 Then the two inequalities can be merged to form the following inequality:

 $x\le2y\le4x\le120$

 Additionally, we must ensure that $2y<60$

 Therefore we must find pairs $(x,y)$ that satisfy the inequality above. A bit of trial and error and observing patterns leads to the answer $480$.

 It should be noted that the cases for $x\le15$ and $x>15$ should be considered separately in order to ensure that $2y < 60$.

 





\textbf{Solution 4 - Unrigorous engineers induction solution}

We will try out small cases.

 By replacing 60 in this problem with 2, we count only 1 ordered pair. By doing with 4, we count 4 ordered pairs. With 6, we get 7 pairs. With 8 we get 12. By continuing on, and then finding the difference between adjacent terms (1,3,3,5,5,...). We suspect that if 60 was replaced with 2n, we will find 1+3+3+5+5+7+7 ...., where there will be n terms. Thus, our answer is 1+3+3+5+5.... 29+29+31 = 16*30 = 480.

 -Alexlikemath

 


 





\textbf{Solution 5 - Counting Head on}

Notice $x$ and $y$ both must be equal or less than 30. The inequalities given have no complicated qualities. We can recompile them by understanding: Two times the larger integer will also be larger than the smaller integer; Two times the smaller integer is greater or equal to the greater integer if and only if the greater integer is less or equal to the double of the smaller integer. Knowing this, we create a chart. We will first solve without order, then multiply pairs by 2 at the end. $x$ can be 1-30, so we'll start with 1. The only possible value for $y$ is 1. For $x = 2$, $y$ can be 2-4. For $x = 3$, $y$ can be 3-6. There is an obvious pattern here. For every integer after 1, the possible values for $y$ will be numbers $x$-$2x$. This predictably ends at $x = 15$ because $y$ will reach 30. When $x = 16$, then the number of possible values of $y$ will begin to drop again, equaling the amount when $x = 14$. Then when we finally sum the group together, for $x = 2$ to $x = 14$ there are 104 pairs with 2 distinct values, and 13 values with congruent values. These will not be multiplied by 2 later on. $x = 16$ to $x = 28$ gives the same amount. Then $x = 15$ and $x = 29$ gives 16 and 2 values respectively, with two congruent values each. Finally, $x = 1$ and $x = 30$ give 3 and 1 respectively.

 Sum them together and you will get 480.

 
-jackshi2006

 




	\item (\textit{AIME I 1998}) The graph of $y^2 + 2xy + 40|x|= 400$ partitions the plane into several regions.  What is the area of the bounded region?


 
 \textbf{Answer: }800



\textbf{Solution}

The equation given can be rewritten as:

 
            \begin{enumerate}[label=(\roman*)]
                \item $40|x| = - y^2 - 2xy + 400$
            \end{enumerate}
            We can split the equation into a piecewise equation by breaking up the \href{https://artofproblemsolving.com/wiki/index.php?title=Absolute\_value}{absolute value}:

 
            \begin{enumerate}[label=(\roman*)]
                \item $40x = -y^2 - 2xy + 400\quad\quad x\ge 0$\item $40x = y^2 + 2xy - 400 \quad \quad x < 0$
            \end{enumerate}
            Factoring the first one: (alternatively, it is also possible to \href{https://artofproblemsolving.com/wiki/index.php?title=Completing\_the\_square}{complete the square})

 
            \begin{enumerate}[label=(\roman*)]
                \item $40x + 2xy  = -y^2 + 400$\item $2x(20 + y)=  (20 - y)(20 + y)$
            \end{enumerate}
            
                    \begin{figure}[H]
                        \centering
                        \includegraphics[width=0.5\textwidth]{Solution/AIME_I_1998_3_Solution_Figure_1.png}
                    \end{figure}
                    Hence, either $y = -20$, or $2x = 20 - y \Longrightarrow y = -2x + 20$.

 Similarily, for the second one, we get $y = 20$ or $y = -2x - 20$. If we graph these four equations, we see that we get a parallelogram with base 20 and height 40. Hence the answer is $\boxed{800}$.

 





\textbf{Solution 2}

The equation can be rewritten as: $(x+y)^2=(|x|-20)^2$. Do casework as above.

 




	\item (\textit{AIME I 1998}) Nine tiles are numbered $1, 2, 3, \cdots, 9,$ respectively.  Each of three players randomly selects and keeps three of the tiles, and sums those three values.  The \href{https://artofproblemsolving.com/wiki/index.php?title=Probability}{probability} that all three players obtain an \href{https://artofproblemsolving.com/wiki/index.php?title=Odd}{odd} sum is $m/n,$ where $m$ and $n$ are \href{https://artofproblemsolving.com/wiki/index.php?title=Relatively\_prime}{relatively prime} \href{https://artofproblemsolving.com/wiki/index.php?title=Positive\_integer}{positive integers}.  Find $m+n.$


 
 \textbf{Answer: }017



\textbf{Solution}

In order for a player to have an odd sum, he must have an odd number of odd tiles: that is, he can either have three odd tiles, or two even tiles and an odd tile.  Thus, since there are $5$ odd tiles and $4$ even tiles, the only possibility is that one player gets $3$ odd tiles and the other two players get $2$ even tiles and $1$ odd tile.  We count the number of ways this can happen. (We will count assuming that it matters in what order the people pick the tiles; the final answer is the same if we assume the opposite, that order doesn't matter.)

 $\dbinom{5}{3} = 10$ choices for the tiles that he gets.  The other two odd tiles can be distributed to the other two players in $2$ ways, and the even tiles can be distributed between them in $\dbinom{4}{2} \cdot \dbinom{2}{2} = 6$ ways.  This gives us a total of  $10 \cdot 2 \cdot 6 = 120$ possibilities in which all three people get odd sums.

 In order to calculate the probability, we need to know the total number of possible distributions for the tiles.  The first player needs three tiles which we can give him in $\dbinom{9}{3} = 84$ ways, and the second player needs three of the remaining six, which we can give him in $\dbinom{6}{3} = 20$ ways. Finally, the third player will simply take the remaining tiles in $1$ way. So, there are $\dbinom{9}{3} \cdot \dbinom{6}{3} \cdot 1 = 84 \cdot 20 = 1680$ ways total to distribute the tiles.

 We must multiply the probability by 3, since any of the 3 players can have the 3 odd tiles.Thus, the total probability is $\frac{360}{1680} = \frac{3}{14},$ so the answer is $3 + 14 = \boxed{017}$.

 





\textbf{Solution 2}

Let $O$ stand for an odd number and $E$ an even.Therefore, one person must pick $OOO$, the other person must pick $OEE$ and the last person must pick $OEE$. Since any permutation of the order of who is picking or change in the order of the even numbers (e.g. $EOE$ instead of $OEE$) doesn't change the probability, we just need to multiply the probability of one case by $\binom{3}{2}^3 = 27$ as there are 27 such cases (by cases I mean ordered triples of ordered multisets $A, B, C$ such that one of them has 3 $O$'s and the other two have two $E$'s and an $O$ in them, respectively.) . Let's do the case $OOO$, $OEE$, $OEE$.$\frac{5}{9} \times \frac{4}{8} \times \frac{3}{7} \times \frac{2}{6} \times \frac{4}{5} \times \frac{3}{4} \times \frac{1}{3} \times \frac{2}{2} \times \frac{1}{1} = \frac{1}{126}$. We now multiply by 27 to get $\frac{3}{14} \implies m + n = \boxed{017}$- whatRthose

 





\textbf{Solution 3}

For this problem, let's think about parity. There are $5$ odd numbers from $1-9$ and there are four even numbers from $1-9$. Since this problem is asking for the probability that the each player gets an odd sum, we also have to calculate the total numeber of ways. 

 In this case, there are only two ways to get an odd sum. Either have the sequence $OOO$ or $OEE$ where the letters $O$ and $E$ stand for odd and even respectively. 

 Since we constrained to only $5$ odds, the only way to do the pairing is 
 \[OOO-OEE-OEE\]There are of couse three ways to choose who gets the three odds. 

 Once we have chosen who has gotten the three odds, we can actually reorder the sequence like this:
 \[OEE-OEE-OOO\]Now since we are choosing in groups, we can ignor order of these terms.

 For the first $OEE$, there are $5$ ways to choose which odd to use, and $\dbinom{4}{2}$ or $6$ ways to choose the evens.

 Great, let's move on to the second $OEE$. There are $4$ odds left, so there are $4$ ways choose the odds. Since there are only two even numbers, they have to go here, so there is only $1$ way to choose the evens. 

 The three odds will now fall in place. We can now multiply all the numbers since they are independent, and we have $5*6*4*3$ or $360$. 

 Since this is a probability question, we have to ask ourselves how many ways are there to distribute the 9 tiles equally among 3 players. 

 Fortunately for us, this is not hard as the first player has $\dbinom{9}{3}$ options and the second player has $\dbinom{6}{3}$. When we multiply these, we get $1680$. This is our denominator.

 When we make the fraction, we have $\frac{360}{1680}$. When we simplify it, we have:
 \[\frac{3}{14}\]We are asked to find the sum of the numerator and the denominator, so summing these, we have:
 \[\boxed{017}\]-Pi\_is\_3.14

 





\textbf{Solution 4}

All three sums are odd if and only if two players select $OEE$ and the other selects $OOO$. The number of possible scenarios for the three players (ignoring how each player chooses three tiles) is $3\times \dbinom{5}{3, 1, 1}\times\dbinom{4}{2, 2}$. The total number of configurations is $\dbinom{9}{3, 3, 3}$, and $\frac{3\times\dbinom{5}{3, 1, 1}\times\dbinom{4}{2, 2}}{\dbinom{9}{3, 3, 3}}=\frac{3}{14}$. The answer is thus 
 \[\boxed{017}\]-cretinouscretin

 




	\item (\textit{AIME I 1998}) Given that $A_k = \frac {k(k - 1)}2\cos\frac {k(k - 1)\pi}2,$ find $\left\lvert A_{19} + A_{20} + \cdots + A_{98}\right\rvert.$


 
 \textbf{Answer: }040



\textbf{Solution}

Though the problem may appear to be quite daunting, it is actually not that difficult. $\frac {k(k-1)}2$ always evaluates to an integer (\href{https://artofproblemsolving.com/wiki/index.php?title=Triangular\_number}{triangular number}), and the \href{https://artofproblemsolving.com/wiki/index.php?title=Cosine}{cosine} of $n\pi$ where $n \in \mathbb{Z}$ is 1 if $n$ is even and -1 if $n$ is odd. $\frac {k(k-1)}2$ will be even if $4|k$ or $4|k-1$, and odd otherwise.  

 So our sum looks something like:

 $\left|\sum_{i=19}^{98} A_i\right| =- \frac{19(18)}{2} + \frac{20(19)}{2} + \frac{21(20)}{2} - \frac{22(21)}{2} - \frac{23(22)}{2} + \frac{24(23)}{2} \cdots - \frac{98(97)}{2}$

 If we group the terms in pairs, we see that we need a formula for $-\frac{(n)(n-1)}{2} + \frac{(n+1)(n)}{2} = \left(\frac n2\right)(n+1 - (n-1)) = n$. So the first two fractions add up to $19$, the next two to $-21$, and so forth.

 If we pair the terms again now, each pair adds up to $-2$. There are $\frac{98-19+1}{2 \cdot 2} = 20$ such pairs, so our answer is $|-2 \cdot 20| = \boxed{040}$.

 




	\item (\textit{AIME I 1998}) Let $ABCD$ be a \href{https://artofproblemsolving.com/wiki/index.php?title=Parallelogram}{parallelogram}.  Extend $\overline{DA}$ through $A$ to a point $P,$ and let $\overline{PC}$ meet $\overline{AB}$ at $Q$ and $\overline{DB}$ at $R.$  Given that $PQ = 735$ and $QR = 112,$ find $RC.$


 
 \textbf{Answer: }308



\textbf{Solution 1}


                    \begin{figure}[H]
                        \centering
                        \includegraphics[width=0.5\textwidth]{Solution/AIME_I_1998_6_Solution_Figure_1.png}
                    \end{figure}
                    

 There are several \href{https://artofproblemsolving.com/wiki/index.php?title=Similar\_triangles}{similar triangles}. $\triangle PAQ\sim \triangle PDC$, so we can write the \href{https://artofproblemsolving.com/wiki/index.php?title=Proportion}{proportion}:

 $\frac{AQ}{CD} = \frac{PQ}{PC} = \frac{735}{112 + 735 + RC} = \frac{735}{847 + RC}$

 Also, $\triangle BRQ\sim DRC$, so:

 $\frac{QR}{RC} = \frac{QB}{CD} = \frac{112}{RC} = \frac{CD - AQ}{CD} = 1 - \frac{AQ}{CD}$


 $\frac{AQ}{CD} = 1 - \frac{112}{RC} = \frac{RC - 112}{RC}$

 Substituting,

 $\frac{AQ}{CD} = \frac{735}{847 + RC} = \frac{RC - 112}{RC}$


 $735RC = (RC + 847)(RC - 112)$
$0 = RC^2 - 112\cdot847$

 Thus, $RC = \sqrt{112*847} = 308$.

 





\textbf{Solution 2}

We have $\triangle BRQ\sim \triangle DRC$ so $\frac{112}{RC} = \frac{BR}{DR}$. We also have $\triangle BRC \sim \triangle DRP$ so $\frac{ RC}{847} = \frac {BR}{DR}$. Equating the two results gives $\frac{112}{RC} =  \frac{ RC}{847}$ and so $RC^2=112*847$ which solves to $RC=\boxed{308}$

 




	\item (\textit{AIME I 1998}) Let $n$ be the number of ordered quadruples $(x_1,x_2,x_3,x_4)$ of positive odd \href{https://artofproblemsolving.com/wiki/index.php?title=Integer}{integers} that satisfy $\sum_{i = 1}^4 x_i = 98.$  Find $\frac n{100}.$


 
 \textbf{Answer: }196



\textbf{Solution 1 (Clever Stars and Bars Manipulation)}

We want $x_1 +x_2+x_3+x_4 =98$. This seems like it can be solved with stars and bars, however note that the quadruples all need to be odd. This motivates us to set $x_i= 2y_i +1$, as for all integers $y_i$, $2y_i + 1$ will be odd. Substituting we get 
 \[2y_1+2y_2+2y_3+2y_4 +4 = 98 \implies y_1+y_2+y_3+y_4 =47\]Note that this is an algebraic bijection, we have simplified the problem and essentially removed the odd condition, so now we can finish with plain stars as bars, which gives us $n= {50\choose3}$. Computing this and dividing by 100 gives us an answer of $\boxed{196}$.~smartguy888

 





\textbf{Solution 2}

Define $x_i = 2y_i - 1$. Then $2\left(\sum_{i = 1}^4 y_i\right) - 4 = 98$, so $\sum_{i = 1}^4 y_i = 51$. 

 So we want to find four natural numbers that sum up to 51; we can imagine this as trying to split up 51 on the number line into 4 ranges. This is equivalent to trying to place 3 markers on the numbers 1 through 50; thus the answer is $n = {50\choose3} = \frac{50 * 49 * 48}{3 * 2} = 19600$, and $\frac n{100} = \boxed{196}$.

 





\textbf{Solution 3}

Another way we can approach this problem is by imagining a line of 98 stones. We want to place the stones into $4$ boxes so that each box has an odd number of stones. We then proceed by placing one stone in each box to begin with, ensuring that we have a positive number in every box. Now we have $94$ stones left. Because we want an odd number in each box, we pair the stones, creating $47$ sets of $2$. Every time we add a pair to one of the boxes, the number of stones in the box remains odd, because (an odd \#) + (an even \#) = (an odd \#).

 Our problem can now be restated: how many ways are there to partition a line of $47$ stones? We can easily solve this by using $3$ sticks to separate the stones into $4$ groups, and this is the same as arranging a line of $3$ sticks and $47$ stones.  
 \[\frac{50!}{47! \cdot 3!} = 19600\] 
 \[\frac{50 * 49 * 48}{3 * 2} = 19600\] Our answer is therefore $\frac{19600}{100} = \boxed{196}$

 


 





\textbf{Solution 4}

Let $x = a + b$ and $y = c + d$. Then $x + y = 98$, where $x, y$ are positive even integers ranging from $2-98$.

 -When $(x, y) = (2, 96)$, $(a, b) = (1, 1)$ and $(c, d) = (1, 95), (3, 93),...,(95, 1)$. This accounts for $48$ solutions.

 -When $(x, y) = (4, 94)$, $(a, b) = (1, 3), (3, 1)$ and $(c, d) = (1, 93),  (3, 91),...,(93, 1)$. This accounts for $94$ solutions.

 We quickly see that the total number of acceptable ordered pairs $(a, b, c, d)$ is:

 
 \begin{align*} &\mathrel{\phantom{=}} 1 \cdot 48 + 2 \cdot 47 + 3 \cdot 46 + ... + 48 \cdot 1\\ &= (24.5 - 23.5)(24.5 + 23.5) + (24.5 - 22.5)(24.5 + 22.5) + ... + (24.5 + 23.5)(24.5 - 23.5)\\ &= 48(24.5)^2 - 2(0.5^2 + 1.5^2 + ... + 23.5^2)\\ &= 28812 - \frac{1^2 + 3^2 + ... + 47^2}{2}\\ &= 28812 - \frac{1^2 + 2^2 + ... + 47^2 - 4(1^2 + 2^2 + ... + 23^2)}{2}\\ &= 28812 - \frac{\frac{47(47 + 1)(2(47) + 1)}{6} - \frac{4(23)(23 + 1)(2(23) + 1)}{6}}{2}\\ &= 19600 \end{align*}


 Therefore, $\frac{n}{100} = \frac{19600}{100} = \boxed{196}$.

 (This solution uses the sum of squares identity to calculate $1^2 + 2^2 + ... + 47^2$ and $1^2 + 2^2 + ... + 23^2$.)

 -baker77

 





\textbf{Solution 5}

We write the generating functions for each of the terms, and obtain $(x+x^3+x^5\cdots)^4$ as the generating function for the sum of the $4$ numbers. We seek the $x^{98}$ coefficient, or the $x^{94}$ coefficient in $(1+x^2+x^4...)^4.$ Now we simplify this as $\left(\frac{1}{1-x^2}\right)^4=\binom{3}{3} +\binom{4}{3}x^2+\binom{5}{3}x^4 \cdots$ and in general we want that the coefficient of $x^{2k}$ is $\binom{k+3}{3}.$ We see the $x^{94}$ coefficient so we let $k=47$ and so the coefficient is $\binom{50}{3}=19600$ in which $\frac{n}{100}=\boxed{196}.$

 




	\item (\textit{AIME I 1998}) Except for the first two terms, each term of the sequence $1000, x, 1000 - x,\ldots$ is obtained by subtracting the preceding term from the one before that.  The last term of the sequence is the first \href{https://artofproblemsolving.com/wiki/index.php?title=Negative}{negative} term encountered.  What positive integer $x$ produces a sequence of maximum length?


 
 \textbf{Answer: }618



\textbf{Solution}

The best way to start is to just write out some terms.

 




	\item (\textit{AIME I 1998}) Two mathematicians take a morning coffee break each day.  They arrive at the cafeteria independently, at random times between 9 a.m. and 10 a.m., and stay for exactly $m$ minutes.  The \href{https://artofproblemsolving.com/wiki/index.php?title=Probability}{probability} that either one arrives while the other is in the cafeteria is $40 \%,$ and $m = a - b\sqrt {c},$ where $a, b,$ and $c$ are \href{https://artofproblemsolving.com/wiki/index.php?title=Positive}{positive} \href{https://artofproblemsolving.com/wiki/index.php?title=Integer}{integers}, and $c$ is not divisible by the square of any \href{https://artofproblemsolving.com/wiki/index.php?title=Prime}{prime}.  Find $a + b + c.$


 
 \textbf{Answer: }087



\textbf{Solution 1}

Let the two mathematicians be $M_1$ and $M_2$.  Consider plotting the times that they are on break on a \href{https://artofproblemsolving.com/wiki/index.php?title=Coordinate\_plane}{coordinate plane} with one axis being the time $M_1$ arrives and the second axis being the time $M_2$ arrives (in minutes past 9 a.m.). The two mathematicians meet each other when $|M_1-M_2| \leq m$. Also because the mathematicians arrive between 9 and 10, $0 \leq M_1,M_2 \leq 60$. Therefore, $60\times 60$ square represents the possible arrival times of the mathematicians, while the shaded region represents the arrival times where they meet.
                    \begin{figure}[H]
                        \centering
                        \includegraphics[width=0.5\textwidth]{Solution/AIME_I_1998_9_Solution_Figure_1.png}
                    \end{figure}
                    It's easier to compute the area of the unshaded region over the area of the total region, which is the probability that the mathematicians do not meet:

 $\frac{(60-m)^2}{60^2} = 0.6$
$(60-m)^2 = 36\cdot 60$
$60 - m = 12\sqrt{15}$
$\Rightarrow m = 60-12\sqrt{15}$

 So the answer is $60 + 12 + 15 = 087$.

 





\textbf{Solution 2}

Case 1:

 
                    \begin{figure}[H]
                        \centering
                        \includegraphics[width=0.5\textwidth]{Solution/AIME_I_1998_9_Solution_Figure_2.png}
                    \end{figure}
                    Case 2:

 
                    \begin{figure}[H]
                        \centering
                        \includegraphics[width=0.5\textwidth]{Solution/AIME_I_1998_9_Solution_Figure_3.png}
                    \end{figure}
                    

 We draw a \href{https://artofproblemsolving.com/wiki/index.php?title=Number\_line}{number line} representing the time interval. If mathematician $M_1$ comes in at the center of the time period, then the two mathematicions will meet if $M_2$ comes in somewhere between $m$ minutes before and after $M_1$ comes (a total range of $2m$ minutes). However, if $M_1$ comes into the cafeteria in the first or last $m$ minutes, then the range in which $M_2$ is reduced to somewhere in between $m$ and $2m$.

 We now try to find the \href{https://artofproblemsolving.com/wiki/index.php?title=Weighted\_average}{weighted average} of the chance that the two meet. In the central $60-2m$ minutes, $M_1$ and $M_2$ have to enter the cafeteria within $m$ minutes of each other; so if we fix point $M_1$ then $M_2$ has a $\frac{2m}{60} = \frac{m}{30}$ probability of meeting.

 In the first and last $2m$ minutes, the probability that the two meet ranges from $\frac{m}{60}$ to $\frac{2m}{60}$, depending upon the location of $M_1$ with respect to the endpoints. Intuitively, the average probability will occur at $\frac{\frac{3}{2}m}{60} = \frac{m}{40}$.

 So the weighted average is:

 
            \begin{enumerate}[label=(\roman*)]
                \item $\frac{\frac{m}{30}(60-2m) + \frac{m}{40}(2m)}{60} = \frac{40}{100}$\item $0 = \frac{m^2}{60} - 2m + 24$\item $0 = m^2 - 120m + 1440$
            \end{enumerate}
            Solving this \href{https://artofproblemsolving.com/wiki/index.php?title=Quadratic\_equation}{quadratic}, we get two roots, $60 \pm 12\sqrt{15}$. However, $m < 60$, so we discard the greater root; and thus our answer $60 + 12 + 15 = 087$.

 \textbf{Video Solution} 




	\item (\textit{AIME I 1998}) Eight \href{https://artofproblemsolving.com/wiki/index.php?title=Sphere}{spheres} of \href{https://artofproblemsolving.com/wiki/index.php?title=Radius}{radius} 100 are placed on a flat \href{https://artofproblemsolving.com/wiki/index.php?title=Plane}{surface} so that each sphere is \href{https://artofproblemsolving.com/wiki/index.php?title=Tangent}{tangent} to two others and their \href{https://artofproblemsolving.com/wiki/index.php?title=Center}{centers} are the vertices of a regular \href{https://artofproblemsolving.com/wiki/index.php?title=Octagon}{octagon}.  A ninth sphere is placed on the flat surface so that it is tangent to each of the other eight spheres.  The radius of this last sphere is $a +b\sqrt {c},$ where $a, b,$ and $c$ are \href{https://artofproblemsolving.com/wiki/index.php?title=Positive}{positive} \href{https://artofproblemsolving.com/wiki/index.php?title=Integer}{integers}, and $c$ is not divisible by the square of any \href{https://artofproblemsolving.com/wiki/index.php?title=Prime}{prime}.  Find $a + b + c$.


 
 \textbf{Answer: }152



\textbf{Solution 1}

The key is to realize the significance that the figures are spheres, not \href{https://artofproblemsolving.com/wiki/index.php?title=Circle}{circles}. The 2D analogue of the diagram onto the flat surface will not contain 8 circles tangent to a ninth one; instead the circles will overlap since the middle sphere has a larger radius and will sort of “bulge” out.

 
                    \begin{figure}[H]
                        \centering
                        \includegraphics[width=0.5\textwidth]{Solution/AIME_I_1998_10_Solution_Figure_1.png}
                    \end{figure}
                    

 Let us examine the relation between one of the outside 8 spheres and the center one (with radius $r$): 

 
                    \begin{figure}[H]
                        \centering
                        \includegraphics[width=0.5\textwidth]{Solution/AIME_I_1998_10_Solution_Figure_2.png}
                    \end{figure}
                    

 If we draw the segment containing the centers and the radii \href{https://artofproblemsolving.com/wiki/index.php?title=Perpendicular}{perpendicular} to the flat surface, we get a \href{https://artofproblemsolving.com/wiki/index.php?title=Trapezoid}{trapezoid}; if we draw the segment parallel to the surface that connects the center of the smaller sphere to the radii of the larger, we get a right triangle. Call that segment $x$. Then by the \href{https://artofproblemsolving.com/wiki/index.php?title=Pythagorean\_Theorem}{Pythagorean Theorem}:

 
 \[x^2 + (r-100)^2 = (r+100)^2 \Longrightarrow x = 20\sqrt{r}\]

 
                    \begin{figure}[H]
                        \centering
                        \includegraphics[width=0.5\textwidth]{Solution/AIME_I_1998_10_Solution_Figure_3.png}
                    \end{figure}
                    

 $x$ is the distance from one of the vertices of the octagon to the center, so the diagonal of the octagon is of length $2x =40\sqrt{r}$. We can draw another \href{https://artofproblemsolving.com/wiki/index.php?title=Right\_triangle}{right triangle} as shown above. One leg has a length of $200$. The other can be found by partitioning the leg into three sections and using $45-45-90 \triangle$s to see that the leg is $100\sqrt{2} + 200 + 100\sqrt{2} = 200(\sqrt{2} + 1)$. Pythagorean Theorem:

 
 \begin{eqnarray*} (40\sqrt{r})^2 &=& 200^2 + [200(\sqrt{2}+1)]^2\\ 1600r &=& 200^2[(1 + \sqrt{2})^2 + 1] \\ r &=& 100 + 50\sqrt{2} \end{eqnarray*}

 Thus $a + b + c = 100 + 50 + 2 = \boxed{152}$.

 





\textbf{Solution 2}

Isolate a triangle, with base length $200$ (a side of the octagon). This triangle is obviously isoceles. Denote the other side length as $x$. Since the interior angle is $45$ degrees (due to the shape being an octagon), then we can apply Law of Cosines to this triangle. We get: 
 \begin{eqnarray*} 200^2 &=& 2x^2 - 2x^2*cos(45^\circ) \\ &=& 2x^2 - 2x^2*\frac{\sqrt{2}}{2} \\ &=& (2-\sqrt{2})x^2 \end{eqnarray*}

 And thus 
 \[x = \frac{200}{\sqrt{2-\sqrt{2}}}\]

 From the above, $x = 20\sqrt{r}$, so we get

 
 \begin{eqnarray*} r &=& (\frac{200}{20(\sqrt{2-\sqrt{2}})})^2 \\ &=& (\frac{10}{\sqrt{2-\sqrt{2}})})^2 \cdot \frac{2+\sqrt{2}}{2+\sqrt{2}} \\ &=& \frac{200 + 100\sqrt{2}}{2} \\ &=& 100 + 50\sqrt{2} \end{eqnarray*}

 And hence the answer is $100+50+2 = \boxed{152}$

 




	\item (\textit{AIME I 1998}) Three of the edges of a \href{https://artofproblemsolving.com/wiki/index.php?title=Cube}{cube} are $\overline{AB}, \overline{BC},$ and $\overline{CD},$ and $\overline{AD}$ is an interior \href{https://artofproblemsolving.com/wiki/index.php?title=Diagonal}{diagonal}.  Points $P, Q,$ and $R$ are on $\overline{AB}, \overline{BC},$ and $\overline{CD},$ respectively, so that $AP = 5, PB = 15, BQ = 15,$ and $CR = 10.$  What is the \href{https://artofproblemsolving.com/wiki/index.php?title=Area}{area} of the \href{https://artofproblemsolving.com/wiki/index.php?title=Polygon}{polygon} that is the \href{https://artofproblemsolving.com/wiki/index.php?title=Intersection}{intersection} of \href{https://artofproblemsolving.com/wiki/index.php?title=Plane}{plane} $PQR$ and the cube?


 
 \textbf{Answer: }525



\textbf{Solution 1}

"For non-asymptote version of image, see \href{https://artofproblemsolving.com/wiki/index.php?title=File:1998\_AIME-11.png}{Image:1998\_AIME-11.png}"

 
                    \begin{figure}[H]
                        \centering
                        \includegraphics[width=0.5\textwidth]{Solution/AIME_I_1998_11_Solution_Figure_1.png}
                    \end{figure}
                    This approach uses \href{https://artofproblemsolving.com/wiki/index.php?title=Analytic\_Geometry}{analytic geometry}. Let $A$ be at the origin, $B$ at $(20,0,0)$, $C$ at $(20,0,20)$, and $D$ at $(20,20,20)$. Thus, $P$ is at $(5,0,0)$, $Q$ is at $(20,0,15)$, and $R$ is at $(20,10,20)$. 

 Let the plane $PQR$ have the equation $ax + by + cz = d$. Using point $P$, we get that $5a = d$. Using point $Q$, we get 
 \[20a + 15c = d \Longrightarrow 4d + 15c = d \Longrightarrow d = -5c\]Using point $R$, we get 
 \[20a + 10b + 20c = d \Longrightarrow 4d + 10b - 4d = d \Longrightarrow d = 10b\]Thus plane $PQR$’s \href{https://artofproblemsolving.com/wiki/index.php?title=Equation}{equation} reduces to 
 \[\frac{d}{5}x + \frac{d}{10}y - \frac{d}{5}z = d \Longrightarrow 2x + y - 2z = 10\]

 We know need to find the intersection of this plane with that of $z = 0$, $z = 20$, $x = 0$, and $y = 20$. After doing a little bit of algebra, the intersections are the lines 
 \[y = -2x + 10\] 
 \[y = -2x + 50\] 
 \[y = 2z + 10\] 
 \[z = x + 5\]Thus, there are three more vertices on the polygon, which are at $(0,10,0)(0,20,5)(15,20,20)$. 

 We can find the lengths of the sides of the polygons now. There are 4 \href{https://artofproblemsolving.com/wiki/index.php?title=Right\_triangle}{right triangles} with legs of length 5 and 10, so their \href{https://artofproblemsolving.com/wiki/index.php?title=Hypotenuse}{hypotenuses} are $5\sqrt{5}$. The other two are of $45-45-90 \triangle$s with legs of length 15, so their hypotenuses are $15\sqrt{2}$. So we have a \href{https://artofproblemsolving.com/wiki/index.php?title=Hexagon}{hexagon} with sides $15\sqrt{2},5\sqrt{5}, 5\sqrt{5},15\sqrt{2}, 5\sqrt{5},5\sqrt{5}$ By \href{https://artofproblemsolving.com/wiki/index.php?title=Symmetry}{symmetry}, we know that opposite angles of the polygon are congruent. We can also calculate the length of the long diagonal by noting that it is of the same length of a face diagonal, making it $20\sqrt{2}$.

 
                    \begin{figure}[H]
                        \centering
                        \includegraphics[width=0.5\textwidth]{Solution/AIME_I_1998_11_Solution_Figure_2.png}
                    \end{figure}
                    The height of the triangles at the top/bottom is $\frac{20\sqrt{2} - 15\sqrt{2}}{2} = \frac{5}{2}\sqrt{2}$. The \href{https://artofproblemsolving.com/wiki/index.php?title=Pythagorean\_Theorem}{Pythagorean Theorem} gives that half of the base of the triangles is $\frac{15}{\sqrt{2}}$. We find that the middle \href{https://artofproblemsolving.com/wiki/index.php?title=Rectangle}{rectangle} is actually a \href{https://artofproblemsolving.com/wiki/index.php?title=Square}{square}, so the total area is $(15\sqrt{2})^2 + 4\left(\frac 12\right)\left(\frac 52\sqrt{2}\right)\left(\frac{15}{\sqrt{2}}\right) = 525$.

 





\textbf{Solution 2}

First, note that whenever the plane intersects two opposite faces of the cube, the resulting line segments must be parallel.  Because they are part of parallel planes (the faces), they must be either parallel or \href{https://artofproblemsolving.com/wiki/index.php?title=Skew}{skew}; they are both part of plane $PQR$, so they cannot be skew.  Therefore, they are \href{https://artofproblemsolving.com/wiki/index.php?title=Parallel}{parallel}.

 Let the cube's vertices be $A$, $B$, $C$, $D$, $E$, $F$, $G$, and $H$, with $A$, $B$, and $C$ on the bottom face as before, $H$ being the other bottom vertex, $D$ directly above $C$, $E$ above $B$, $F$ above $A$, and $G$ above $H$.

 Clearly, the next vertex of the intersection (starting with $P$, $Q$, $R$) will be somewhere on $DG$.  Let it be $X$, and have a distance of $x$ from D, and a distance of $20 - x$ from $G$.

 Then, the next vertex will be somewhere on $FG$.  It must be parallel to $PQ$, so this implies that it has a distance of $20 - x$ from $G$, and thus a distance of $x$ from $F$.

 Now, the next vertex (call it $Y$) will be somewhere on $AF$.  The segment must be parallel to $QR$, so $FY$ must have length $2x$,  and $AY$ must be $20 - 2x$.

 Since $DX \parallel AP$, $DR \parallel AY$, and $RX \parallel PY$, we must have $\triangle{APY} \sim \triangle{DXR}$; therefore, 
 \[\frac{AP}{DX}=\frac{AY}{DR}\] 
 \[\frac{5}{x}=\frac{20-2x}{10}\] 
 \[x^{2}-10x+25=0\] 
 \[x=5\]

 We can now find that the hexagon has side lengths $15\sqrt {2}$, $5\sqrt {5}$, $5\sqrt {5}$, $15\sqrt {2}$, $5\sqrt {5}$, and $5\sqrt {5}$.  Moreover, opposite angles of this must be equal (by symmetry), so segment $RY$ divides the \href{https://artofproblemsolving.com/wiki/index.php?title=Hexagon}{hexagon} into two \href{https://artofproblemsolving.com/wiki/index.php?title=Isosceles\_trapezoid}{isosceles trapezoids}.  It is easy to find the length of $RY$ (they're midpoints of opposite edges, so the distance between the two points is equal to a face diagonal of the cube, or $20\sqrt {2}$), so it is now easy to finish the problem. From here, we can continue as in the first solution.

 




	\item (\textit{AIME I 1998}) Let $ABC$ be \href{https://artofproblemsolving.com/wiki/index.php?title=Equilateral\_triangle}{equilateral}, and $D, E,$ and $F$ be the \href{https://artofproblemsolving.com/wiki/index.php?title=Midpoint}{midpoints} of $\overline{BC}, \overline{CA},$ and $\overline{AB},$ respectively.  There exist \href{https://artofproblemsolving.com/wiki/index.php?title=Point}{points} $P, Q,$ and $R$ on $\overline{DE}, \overline{EF},$ and $\overline{FD},$ respectively, with the property that $P$ is on $\overline{CQ}, Q$ is on $\overline{AR},$ and $R$ is on $\overline{BP}.$  The \href{https://artofproblemsolving.com/wiki/index.php?title=Ratio}{ratio} of the area of triangle $ABC$ to the area of triangle $PQR$ is $a + b\sqrt {c},$ where $a, b,$ and $c$ are integers, and $c$ is not divisible by the square of any \href{https://artofproblemsolving.com/wiki/index.php?title=Prime}{prime}.  What is $a^{2} + b^{2} + c^{2}$?


 
 \textbf{Answer: }083



\textbf{Solution 1}

The diagram below shows an incorrect version of the problem. To accurately follow the diagram, the problem should read:

 Let $ABC$ be \href{https://artofproblemsolving.com/wiki/index.php?title=Equilateral\_triangle}{equilateral}, and $D, E,$ and $F$ be the \href{https://artofproblemsolving.com/wiki/index.php?title=Midpoint}{midpoints} of $\overline{CA}, \overline{AB},$ and $\overline{BC},$ respectively.  There exist \href{https://artofproblemsolving.com/wiki/index.php?title=Point}{points} $P, Q,$ and $R$ on $\overline{DE}, \overline{EF},$ and $\overline{FD},$ respectively, with the property that $P$ is on $\overline{AQ}, Q$ is on $\overline{BR},$ and $R$ is on $\overline{CP}.$  The \href{https://artofproblemsolving.com/wiki/index.php?title=Ratio}{ratio} of the area of triangle $ABC$ to the area of triangle $PQR$ is $a + b\sqrt {c},$ where $a, b$ and $c$ are integers, and $c$ is not divisible by the square of any \href{https://artofproblemsolving.com/wiki/index.php?title=Prime}{prime}.  What is $a^{2} + b^{2} + c^{2}$?

 


 
                    \begin{figure}[H]
                        \centering
                        \includegraphics[width=0.5\textwidth]{Solution/AIME_I_1998_12_Solution_Figure_1.png}
                    \end{figure}
                    

 WLOG, assume that $AB = BC = AC = 2$.

 We let $x = EP = FQ$, $y = EQ$, $k = PQ$. Since $AE = \frac {1}{2}AB$ and $AD = \frac {1}{2}AC$, $\triangle AED \sim \triangle ABC$ and $ED \parallel BC$.

 By alternate interior angles, we have $\angle PEQ = \angle BFQ$ and $\angle EPQ = \angle FBQ$. By vertical angles, $\angle EQP = \angle FQB$. 

 Thus $\triangle EQP \sim \triangle FQB$, so $\frac {EP}{EQ} = \frac {FB}{FQ}\Longrightarrow\frac {x}{y} = \frac {1}{x}\Longrightarrow x^{2} = y$.

 Since $\triangle EDF$ is equilateral, $EQ + FQ = EF = BF = 1\Longrightarrow x + y = 1$. Solving for $x$ and $y$ using $x^{2} = y$ and $x + y = 1$ gives $x = \frac {\sqrt {5} - 1}{2}$ and $y = \frac {3 - \sqrt {5}}{2}$.

 Using the \href{https://artofproblemsolving.com/wiki/index.php?title=Law\_of\_Cosines}{Law of Cosines}, we get

 $k^{2}  =  x^{2} + y^{2} - 2xy\cos{\frac {\pi}{3}}$ 
            \begin{enumerate}[label=(\roman*)]
                \item $=  \left(\frac {\sqrt {5} - 1}{2}\right)^{2} + \left(\frac {3 - \sqrt {5}}{2}\right)^{2} - 2\left(\frac {\sqrt {5} - 1}{2}\right)\left(\frac {3 - \sqrt {5}}{2}\right)\cos{\frac {\pi}{3}}$\item $= 7 - 3\sqrt {5}$
            \end{enumerate}
            We want the ratio of the squares of the sides, so $\frac {(2)^{2}}{k^{2}} = \frac {4}{7 - 3\sqrt {5}} = 7 + 3\sqrt {5}$ so $a^{2} + b^{2} + c^{2} = 7^{2} + 3^{2} + 5^{2} = \boxed{083}$.

 





\textbf{Solution 2}

WLOG, let $\Delta ABC$ have side length $2.$ Then, $DE = EF = FD = 1.$ We also notice that $\angle CEP = \angle DEF = 60^{\circ},$ meaning $\angle CEF = \angle CEP + \angle DEF = 120^{\circ}.$

 Let $EP = x.$ Since $FQ = x$ by congruent triangles $\Delta EPC$ and $\Delta FQA,$ $EQ = EF - FQ = 1-x.$ We can now apply Law of Cosines to $\Delta CEP, \Delta PEQ,$ and $\Delta CEQ.$

 By LoC on $\Delta CEP,$ we get 
 \[CP^2 = 1^2 + x^2 - 2\cdot 1\cdot x\cdot \left(\frac{1}{2}\right) = x^2 - x + 1.\]

 In a similar vein, using LoC on $\Delta PEQ$ and $\Delta CEQ,$ respectively, earns 
 \[PQ^2 = x^2 + (1-x)^2 - 2\cdot x\cdot (1-x)\cdot \left(\frac{1}{2}\right) = 3x^2 - 3x + 1\] 
 \[CQ^2 = 1^2 + (1-x)^2 - 2\cdot 1\cdot (1-x)\cdot \left(-\frac{1}{2}\right) = x^2 - 3x + 3\]

 We have $CP^2, PQ^2,$ and $CQ^2.$ Additionally, by segment addition, $CP + PQ = CQ.$ Solving for $CP, PQ,$ and $CQ$ from the Law of Cosines expressions and plugging them into the segment addition gets the (admittedly-ugly) equation 
 \[\sqrt{x^2-3x+3} = \sqrt{x^2-x+1} + \sqrt{3x^2-3x+1}.\]

 Since the equation is ugly, we look at what the problem is asking for us to solve. We want $\frac{[ABC]}{[PQR]}.$ We see that $[ABC] = \sqrt{3}$ and $[PQR] = [DEF] - [PDR] - [RFQ] - [QEP] = \frac{\sqrt{3}}{4} - \frac{3}{2}\left(\frac{\sqrt{3}}{2}x(1-x)\right),$ since $[PDR] = [RFQ] = [QEP] = \frac{1}{2}x(1-x)\frac{\sqrt{3}}{2}$ from the sine area formula. Simplifying $\frac{[ABC]}{[PQR]}$ gets us wanting to find $\frac{4}{3x^2-3x+1}.$

 We see $3x^2-3x+1$ in both the denominator of what we want and under a radicand in our algebraic expression, which leads us to think the calculations may not be that bad. Isolate $\sqrt{3x^2-3x+1}$ and square to get 
 \[3x^2-3x+1 = 2x^2-4x+4-2\sqrt{(x^2-3x+3)(x^2-x+1)}\]

 Isolate the radicand and square and expand to get $x^4+2x^3-5x^2-6x+9=4x^4-16x^3+28x^2-24x+12,$ and moving terms to one side and dividing by $3,$ we get 
 \[x^4-6x^3+11x^2-6x+1=0.\]

 This can be factored into $(x^2-3x+1)^2 = 0 \rightarrow x^2-3x+1 = 0 \rightarrow x = \frac{3 \pm \sqrt{5}}{2}.$ From the equation $x^2-3x+1=0,$ we have $x^2=3x+1,$ so plugging that value into the expression we want to find, we get $\frac{4}{3(3x+1)-3x+1} = \frac{4}{6x+2}.$

 Substituting $x = \frac{3-\sqrt{5}}{2}$ into $\frac{4}{6x+2}$ gets an expression of $7+3\sqrt{5},$ so $a^2+b^2+c^2 = \boxed{083}$.

 -PureSwag

 




	\item (\textit{AIME I 1998}) If $\{a_1,a_2,a_3,\ldots,a_n\}$ is a \href{https://artofproblemsolving.com/wiki/index.php?title=Set}{set} of \href{https://artofproblemsolving.com/wiki/index.php?title=Real\_numbers}{real numbers}, indexed so that $a_1 < a_2 < a_3 < \cdots < a_n,$ its \textit{complex power sum} is defined to be $a_1i + a_2i^2+ a_3i^3 + \cdots + a_ni^n,$ where $i^2 = - 1.$  Let $S_n$ be the sum of the complex power sums of all nonempty \href{https://artofproblemsolving.com/wiki/index.php?title=Subset}{subsets} of $\{1,2,\ldots,n\}.$  Given that $S_8 = - 176 - 64i$ and $S_9 = p + qi,$ where $p$ and $q$ are integers, find $\left\lvert p\right\rvert + \left\lvert q\right\rvert.$


 
 \textbf{Answer: }368



\textbf{Solution}

We note that the number of subsets (for now, including the empty subset, which we will just define to have a power sum of zero) with $9$ in it is equal to the number of subsets without a $9$. To easily see this, take all possible subsets of $\{1,2,\ldots,8\}$. Since the sets are ordered, a $9$ must go at the end; hence we can just append a $9$ to any of those subsets to get a new one. 

 Now that we have drawn that \href{https://artofproblemsolving.com/wiki/index.php?title=Bijection}{bijection}, we can calculate the complex power sum recursively. Since appending a $9$ to a subset doesn't change anything about that subset's complex power sum besides adding an additional term, we have that $S_9 = 2S_8 + T_9$, where $T_9$ refers to the sum of all of the $9i^x$. 

 If a subset of size 1 has a 9, then its power sum must be $9i$, and there is only $1$ of these such subsets. There are ${8\choose1}$ with $9\cdot i^2$, ${8\choose2}$ with $9\cdot i^3$, and so forth. So $T_9 =\sum_{k=0}^{8} 9{8\choose{k}}i^{k+1}$. This is exactly the \href{https://artofproblemsolving.com/wiki/index.php?title=Binomial\_expansion}{binomial expansion} of $9i \cdot (1+i)^8$. We can use \href{https://artofproblemsolving.com/wiki/index.php?title=De\_Moivre\%27s\_Theorem}{De Moivre's Theorem} to calculate the power: $(\sqrt{2})^8\cos{8\cdot45} = 16$. Hence $T_9 = 16\cdot9i = 144i$, and $S_9 = 2S_8 + 144i = 2(-176 -64i) + 144i = -352 + 16i$. Thus, $\left\lvert p\right\rvert + \left\lvert q\right\rvert = \left\lvert -352\right\rvert + \left\lvert 16\right\rvert = \boxed{368}$.

 


 




	\item (\textit{AIME I 1998}) An $m\times n\times p$ rectangular box has half the volume of an $(m + 2)\times(n + 2)\times(p + 2)$ rectangular box, where $m, n,$ and $p$ are integers, and $m\le n\le p.$  What is the largest possible value of $p$?


 
 \textbf{Answer: }130



\textbf{Solution 1}


 \[2mnp = (m+2)(n+2)(p+2)\]

 Let’s solve for $p$:

 
 \[(2mn)p = p(m+2)(n+2) + 2(m+2)(n+2)\]
 \[[2mn - (m+2)(n+2)]p = 2(m+2)(n+2)\]
 \[p = \frac{2(m+2)(n+2)}{mn - 2n - 2m - 4} = \frac{2(m+2)(n+2)}{(m-2)(n-2) - 8}\]

 Clearly, we want to minimize the denominator, so we test $(m-2)(n-2) - 8 = 1 \Longrightarrow (m-2)(n-2) = 9$. The possible pairs of factors of $9$ are $(1,9)(3,3)$. These give $m = 3, n = 11$ and $m = 5, n = 5$ respectively. Substituting into the numerator, we see that the first pair gives $130$, while the second pair gives $98$. We now check that $130$ is optimal, setting $a=m-2$, $b=n-2$ in order to simplify calculations. Since
 \[0 \le (a-1)(b-1) \implies a+b \le ab+1\]We have
 \[p = \frac{2(a+4)(b+4)}{ab-8} = \frac{2ab+8(a+b)+32}{ab-8} \le \frac{2ab+8(ab+1)+32}{ab-8} = 10 + \frac{120}{ab-8} \le 130\]Where we see $(m,n)=(3,11)$ gives us our maximum value of $\boxed{130}$.

 
                \begin{itemize}
                    \item Note that $0 \le (a-1)(b-1)$ assumes $m,n \ge 3$, but this is clear as $\frac{2m}{m+2} = \frac{(n+2)(p+2)}{np} > 1$ and similarly for $n$.
                \end{itemize}


                





\textbf{Solution 2}

Similarly as above, we solve for $p,$ but we express the denominator differently:

 
 \[p=\dfrac{2(m+2)(n+2)}{(m+2)(n+2)-4(m+n+2)} \implies \dfrac{1}{p}=\dfrac{1}{2}-\dfrac{2(m+n+2)}{(m+2)(n+2)}.\]Hence, it suffices to maximize $\dfrac{m+n+2}{(m+2)(n+2)},$ under the conditions that $p$ is a positive integer.

 Then since $\dfrac{m+n+2}{(m+2)(n+2)}>\dfrac{1}{2}$ for $m=1,2,$ we fix $m=3.$
 \[\implies \dfrac{1}{p}=\dfrac{1}{2}-\dfrac{2(n+5)}{5(n+2)}=\dfrac{n-10}{10(n+2)},\] where we simply let $n=11$ to achieve $p=\boxed{130}.$

 




	\item (\textit{AIME I 1998}) Define a \textit{domino} to be an ordered pair of \textit{distinct} positive integers. A \textit{proper sequence} of dominos is a list of distinct dominos in which the first coordinate of each pair after the first equals the second coordinate of the immediately preceding pair, and in which $(i,j)$ and $(j,i)$ do not \textit{both} appear for any $i$ and $j$.  Let $D_{40}$ be the set of all dominos whose coordinates are no larger than 40.  Find the length of the longest proper sequence of dominos that can be formed using the dominos of $D_{40}.$


 
 \textbf{Answer: }761



\textbf{Solution 1}

We can draw a comparison between the domino a set of 40 points (labeled 1 through 40) in which every point is connected with every other point. The connections represent the dominoes.

 You need to have all even number of \href{https://artofproblemsolving.com/wiki/index.php?title=Segment}{segments} coming from each point except 0 or 2 which have an odd number of segments coming from the point. (Reasoning for this: Everytime you go to a vertex, you have to leave the vertex, so every vertex reached is equivalent to adding 2 more segments. So the degree of each vertex must be even, with the exception of endpoints) Since there are 39 segments coming from each point it is impossible to touch every segment.

 But you can get up to 38 on each segment because you go in to the point then out on a different segment.  Counting going out from the starting and ending at the ending point we have:

 $\frac{38\cdot 38 + 2\cdot 39}2 = \boxed{761}$

 \textbf{Clarification} To clarify the above solution, every time you move to a new vertex, you take one path in and must leave by another path. Therefore every vertex needs to have an even number of segments leaving it (with the exception of the first and last), because the "in" and "out" segments must make a pair.

 





\textbf{Solution 2}

A proper sequence can be represented by writing the common coordinates of adjacent ordered pairs once. For example, represent (4,7),(7,3),(3,5) as $4,7,3,5 .$ Label the vertices of a regular $n$ -gon $1,2,3, \ldots, n .$ Each domino is thereby represented by a directed segment from one vertex of the $n$ -gon to another, and a proper sequence is represented as a path that retraces no segment. Each time that such a path reaches a non-terminal vertex, it must leave it. 

 
Thus, when $n$ is even, it is not possible for such a path to trace every segment, for an odd number of segments emanate from each vertex. By removing $\frac{1}{2}(n-2)$ suitable segments, however, it can be arranged that $n-2$ segments will emanate from $n-2$ of the vertices and that an odd number of segments will emanate from exactly two of the vertices. 

 
In this situation, a path can be found that traces every remaining segment exactly once, starting at one of the two exceptional vertices and finishing at the other. This path will have length $\binom{n}{2}-\frac{1}{2}(n-2),$ which is 761 when $n=40$.

 \textbf{Note 1} When $n$ is odd, a proper sequence of length $\binom{n}{2}$ can be found using the dominos of $D_{n}$. In this case, the second coordinate of the final domino equals the first coordinate of the first domino. 

 





\textbf{Solution 3}

Let $A_{n}=\{1,2,3, \ldots, n\}$ and $D_{n}$ be the set of dominos that can be formed using integers in $A_{n} .$ Each $k$ in $A_{n}$ appears in $2(n-1)$ dominos in $D_{n},$ hence appears at most $n-1$ times in a proper sequence from $D_{n}.$ Except possibly for the integers $i$ and $j$ that begin and end a proper sequence, every integer appears an even number of times in the sequence. 

 
Thus, if $n$ is even, each integer different from $i$ and $j$ appears on at most $n-2$ dominos in the sequence, because $n-2$ is even, and $i$ and $j$ themselves appear on at most $n-1$ dominos each. This gives an upper bound of
 \[\frac{1}{2}\left[(n-2)^{2}+2(n-1)\right]=\frac{n^{2}-2 n+2}{2}\]dominos in the longest proper sequence in $D_{n}.$ This bound is in fact attained for every even $n.$ It is easy to verify this for $n=2$, so assume inductively that a sequence of this length has been found for a particular value of $n$. 

 
Without loss of generality, assume $i=1$ and $j=2,$ and let $_{p} X_{p+2}$ denote a four-domino sequence of the form $(p, n+1)(n+1, p+1)(p+1, n+2)(n+2, p+2) .$ By appending
 \[{ }_{2} X_{4},{ }_{4} X_{6}, \ldots,{ }_{n-2} X_{n},(n, n+1)(n+1,1)(1, n+2)(n+2,2)\]to the given proper sequence, a proper sequence of length
 \[\frac{n^{2}-2 n+2}{2}+4 \cdot \frac{n-2}{2}+4=\frac{n^{2}+2 n+2}{2}=\frac{(n+2)^{2}-2(n+2)+2}{2}\]is obtained that starts at $i=1$ and ends at $j=2 .$ This completes the inductive proof. 

 
In particular, the longest proper sequence when $n=40$ is 761.

 \textbf{Note 2} In the language of graph theory, this is an example of an Eulerian circuit.

 





\textbf{Solution 4}

Consider the segments joining the vertices of a regular $n$-gon. For odd $n$, we see that the number of segments is quite easily $\binom{n-1}{2}$. This is because every vertex touches every other vertex the same number of times. ($\frac{n-1}{2}$ times to be exact). Hence the answer for odd cases is $n\frac{n-1}{2}=\binom{n-1}{2}$. (This is because a segment that starts at the first vertex also ends at the first vertex).For even $n$ however, every vertex touches $\frac{n-2}{2}$ vertices. However, one may be motivated to say that the answer (as in the odd case) is $n\frac{n-2}{2}$. But this is incorrect, because for the even case, it never ends at the vertex (the first vertex) you started at. So it must end at another vertex. But that vertex has already $\frac{n-2}{2}$ other segments touching it. So we have that the final answer is $1$ plus $n\frac{n-2}{2}$. The case for $n=40$, is $761$.

 \textbf{Note} The reason it touches every single other vertex $\frac{n-1}{2}$ for odd $n$ is because there are a total of $\binom{n-1}{2}$ segments, and once dividing $\binom{n-1}{2}$ by $n$, you will then have the number of segments that are connected to each vertex.For the even case every vertex has at least $\frac{n-2}{2}$ other segments touching it. This is because (you CAN convince yourself through a painful induction/observation argument, but you probably shouldn't) $n$ is even, and if you try to apply the odd case to the even case, (namely that there is $\frac{n-1}{2}$ segments touching each vertex), there would have to then be $n\frac{n-1}{2}$ total segments, which never works since it never loops back to the vertex you started on). So at LEAST, there should be $\frac{n-2}{2}$ segments touching each vertex. However, there is also at most $\frac{n-2}{2}$ segments touching each vertex. Hence by the (what I call) the "less-than-or-greater-than" argument, there must be $\frac{n-2}{2}$ segments out of each vertex. (Mind the plus one extra vertex since once again, the vertex you started on, it doesn't loop around, so it must end at another vertex). Hence the answer is $n\frac{n-2}{2}+1$. 

 (I am not very good at explaining things. Sorry if it didn't make sense. Maybe if you find some way to contact me on aops, I could try and help). (For instance, I could see possible confusion at the part where I claim that the minimum lines be $\frac{n-2}{2}$. You might think, "Oh why not $\frac{n-1}{2}-1$, or even $\frac{n-1}{2}-2$ or even subtract $3$?" Well if you actually start off with the fact that there should be at MOST $\frac{n-2}{2}$ segments from each vertex, then it is obvious, since then there must be a minimum of $\frac{n-2}{2}$ segments from each vertex).

 





\textbf{Solution 5}

We can see that $|D_{40}| = 780$, since we are choosing 2 integers $[1,40]$, and order doesn't matter because $(i,j)$ and $(j,i)$ aren't both in the set. Then from doing a smaller example of $D_4$, we can note that non-endpoints must have an even number of pairs in $D$ in order for one domino's end to match another's beginning. Then, in order to maximize the length we want to minimize the number of dominoes we remove to make all pairs be even (except endpoints). Then we can see that for every two pairs with an odd number of pairs, we can connect single ends to form a circular chain (i.e. if we want to ensure that 2,3 have even pairs then we can imagine (3,1)(1,2)). Thus we can divide 40 by 4 to get the number of numbers to remove. Thus we need to remove 10 numbers. But each number is connected to another, so we remove 20 dominoes. Thus there are 780-20 = 760 dominoes not including endpoints. Finally, we can add one more domino at the end since it doesn't need to match up with the first of another to make it 761.

 





\textbf{Solution 6 (Primitive Modulo Number Theory)}

$\textbf{Lemma}$:

 A sequence of $n$ domino's where $n$ is even will satisfy the longest sequence of proper domino's with length $n(\frac{n}{2} -1) + 1$.

 $\textbf{Proof of Lemma}$:

 Every domino falls into a primitive modulo class $n$ residue. Thus, for all pairs $(i,j)$, we shall only take their modulo, and that will represent our proper string. 

 To create the largest set, we must start with the domino $(1,2)$. We then incrementally increase each position by 1, obtaining $(2,3)$, $(3,4)$, and so on until $(n-1,n)$. From here, we add one more to $n$ given the primitive modulo and obtain the new ordered pair $(n,1)$. Our next sequence starts with $(1,3)$, then $(2,4)$, and so on until $(n-2, n)$ is obtained, from which the next term is $(n-1, 1)$, and then $(n,2)$. Our new sequence begins with 2 and increases accordingly: if the last element in the modulo cycle is $j$ from $(i,j)$, the next term will be $(j, 2j+1)$. 

 Now consider the section where the term $(n, \frac{n}{2} - 1)$ is obtained. This occurs in sequence $\frac{n}{2} - 1$. From here, the next term must increase by $\frac{n}{2}$, and we proceed until the term $(\frac{n}{2} - 1, n - 1)$ is obtained. The next term obtained is $(\frac{n}{2}, n)$. However, the next term of this sequence is $(\frac{n}{2} + 1, 1)$. We know that this term appeared before, particularly, when we were adding by $(\frac{n}{2}-1)$. Thus, the string ends here, at the $\frac{n}{2} - 1$ cycle. We have to account for the term before we found out where the string ends, particularly, $(\frac{n}{2}, n)$, and so there are $n(\frac{n}{2} -1) + 1$ elements.

 Thus, the longest string obtained falls into $n(\frac{n}{2} -1) + 1$ for all even $n$. $\square$

 $\textbf{1998 AIME Q15}$

 We apply the Lemma with $n = 40$ to obtain $40 \cdot(20-1) + 1 = \boxed{761}$.

 




\end{enumerate}